%=========================================================================
% Curriculum Vitae -- Zhengxu Yu
% Layout modelled on a classic LaTeX academic CV:
%   centred small-caps name, two-column ruled entries, sans-serif headings.
%=========================================================================
\documentclass[11pt,letterpaper]{article}

\usepackage[T1]{fontenc}
\usepackage[utf8]{inputenc}
\usepackage{charter}           % Bitstream Charter body text
\usepackage[scaled=0.9]{helvet}
\usepackage{array}
\usepackage{longtable}
\usepackage{enumitem}
\usepackage{xcolor}
\usepackage{microtype}
\usepackage[margin=1in,top=0.85in,bottom=0.8in]{geometry}
\usepackage[colorlinks=true,urlcolor=linkcolor,linkcolor=linkcolor,
            citecolor=linkcolor,breaklinks=true]{hyperref}

\definecolor{linkcolor}{RGB}{125,25,45}

%------------------------------------------------- CV variant switch
%   \researchtrue   research CV: research overview first, then education, employment,
%                   full publication list
%   \researchfalse  industry CV: no overview, employment first, then education, trimmed list
\newif\ifresearch
\researchtrue

%------------------------------------------------- global look
\pagestyle{plain}
\setlength{\parindent}{0pt}
\setlength{\LTpre}{0pt}
\setlength{\LTpost}{0pt}
\renewcommand{\arraystretch}{1.05}

%------------------------------------------------- section headings
\makeatletter
\renewcommand{\section}[1]{%
  \vspace{1.05em}%
  {\sffamily\bfseries\small\color{linkcolor}\MakeUppercase{#1}}\par
  \vspace{-0.25em}{\color{linkcolor}\rule{\textwidth}{0.9pt}}\par
  \vspace{0.55em}%
  \@afterindentfalse\@afterheading}
\renewcommand{\subsection}[1]{%
  \vspace{0.7em}%
  {\sffamily\bfseries\small #1}\par
  \vspace{0.35em}%
  \@afterindentfalse\@afterheading}
\makeatother

%------------------------------------------------- ruled two-column entries
% \begin{entries}{<left column width>} ... \end{entries}
\newenvironment{entries}[1]
  {\begin{longtable}{@{}>{\raggedleft\arraybackslash}p{#1}@{\hspace{8pt}}|@{\hspace{8pt}}p{\dimexpr\textwidth-#1-18pt\relax}@{}}}
  {\end{longtable}}

% vertical strut that keeps the rule continuous between entries
\newcommand{\gap}[1][14pt]{\rule{0pt}{#1}}

% right-aligned date column (awards, talks, service)
\newenvironment{datedentries}[1][3cm]
  {\begin{longtable}{@{}p{\dimexpr\textwidth-#1-18pt\relax}@{\hspace{8pt}}|@{\hspace{8pt}}>{\raggedleft\arraybackslash}p{#1}@{}}}
  {\end{longtable}}

% simple unruled list -- use while an entry group has no dates yet
\newenvironment{plainentries}
  {\begin{list}{}{\setlength{\leftmargin}{0pt}\setlength{\itemsep}{2pt}%
     \setlength{\parsep}{0pt}\setlength{\topsep}{0pt}}}
  {\end{list}}

%------------------------------------------------- publication list
\newenvironment{publist}
  {\begin{list}{}{%
     \setlength{\labelwidth}{2.6em}\setlength{\labelsep}{0.5em}%
     \setlength{\leftmargin}{3.1em}\setlength{\itemsep}{0.45em}%
     \setlength{\parsep}{0pt}\setlength{\topsep}{0.4em}%
     \renewcommand{\makelabel}[1]{\hfill\textbf{##1}}}}
  {\end{list}}
\newcommand{\pub}[1]{\item[#1]}
\newcommand{\me}{\underline{Zhengxu Yu}}
\newcommand{\venue}[1]{\textit{#1}}

\begin{document}

%========================================================== header
\begin{center}
  {\Huge\bfseries Zhengxu Yu}
\end{center}
\vspace{0.8em}

\begin{minipage}[t]{0.5\textwidth}
Huawei Technologies R\&D UK\\
London, United Kingdom
\end{minipage}%
\begin{minipage}[t]{0.5\textwidth}
\raggedleft
\texttt{\href{mailto:yuzxfred@gmail.com}{yuzxfred@gmail.com}}\\
\texttt{\href{https://zhengxuyu.github.io/}{https://zhengxuyu.github.io/}}\\
\texttt{\href{https://dblp.org/pid/246/3155}{https://dblp.org/pid/246/3155}}
\end{minipage}

%========================================================== research overview
\ifresearch
\section{Research Overview}

My research concerns the generalisation limits of learning systems: how far competence acquired
from finite experience extends beyond the distribution that produced it. Existing accounts of
distribution shift buy their guarantees with assumptions that cannot be checked at deployment, and
the methods built on them have proved hard to separate from plain empirical risk minimisation. My
interest is in the case they do not reach, where the new regime is not a perturbation of the
training one and no amount of interpolation within what was seen will produce the right behaviour.

What is missing there is exploration that can leave the region the training data describes rather
than search inside it, and I take that to depend on how the agent represents its environment.
Dynamics learned as executable programs, written from the agent's own interaction and revised
wherever prediction and observation disagree, form a hypothesis that can be rejected outright and
replaced in a way a distributed representation does not allow. Whether this is what governs
behaviour in an unfamiliar regime is a question about held-out environments and has to be settled
by running there: currently on ARC-AGI-3, and in self-organising multi-agent populations.
\fi

%========================================================== education / employment
% Defined here, emitted below in an order that depends on the variant: a research CV
% leads with training, an industry CV leads with what was shipped.
\newcommand{\educationsection}{%
\section{Education}

\begin{entries}{1.7in}
\gap[4pt]\textbf{Ph.D.}\par Computer Science\par 2017--2021 &
  \gap[4pt]Zhejiang University\par
  Advisors: Prof.\ Deng Cai and Prof.\ Xiaofei He\par
  Thesis: Stable Learning for Non-I.I.D Data\\
\gap\textbf{M.Sc.}\par Computer Science\par 2015--2016 &
  \gap University of Surrey\par
  Advisor: Prof.\ H.\ Lilian Tang\par
  Thesis: CNN-based Mycobacterium Cells Segmentation for Time-lapse Images \\
\gap\textbf{B.Eng.}\par Electrical Engineering\par 2011--2015 &
  \gap Jilin University \\
\end{entries}
}

\newcommand{\employmentsection}{%
\section{Employment}

\begin{entries}{1.7in}
\gap[4pt]\textbf{Research Scientist}\par Jan 2026--Present &
  \gap[4pt]Huawei Technologies R\&D UK, London\par
  - Built a generative code world model in which agents explore open-ended environments and
  encode the rules they discover as executable, self-improving code, reaching 100\% Relative
  Human Action Efficiency (RHAE) on the ARC-AGI-3 community leaderboard.\par
  - Led research and system design for \textbf{OneManCompany (OMC)}, an open-source agentic AI
  platform for self-organising multi-agent systems, leading a 5-person team; 400+ GitHub
  stars, with the design published as arXiv:2604.22446.\par
  \url{https://github.com/1mancompany/OneManCompany} \\
\gap\textbf{Member of Technical Team}\par Aug 2025--Jan 2026 &
  \gap Inephany Ltd., London\par
  - Built the multi-agent reinforcement learning training infrastructure for large language
  models and vision transformers, supporting distributed training runs on clusters of
  500+ GPUs. \\
\gap\textbf{Algorithm Expert}\par 2024--2025 &
  \gap Apsara Lab (now Tongyi Lab), Alibaba Cloud\par
  - Developed RL-based post-training methods for LLM reasoning efficiency, alignment
  and agentic tool use, cutting reasoning-token usage by 50\% on Qwen-7B at no accuracy loss
  (NeurIPS 2025, AAAI 2026).\par
  - Delivered three bespoke LLM agent systems with dynamic reasoning and tool use for
  decision support and information retrieval, deployed for the International Olympic
  Committee and several municipal government bureaus in China.\par
  - Shipped those three systems and two further large-scale enterprise solutions as tech lead
  and core algorithm owner, contributing to over RMB 25M in cumulative GAAP-recognised
  revenue.\par
  - Led a 13-engineer feature team and a 40-person data-annotation group, supervising 8
  research interns and junior engineers. \\
\gap\textbf{Senior Algorithm Engineer}\par 2021--2024 &
  \gap Apsara Lab (now Tongyi Lab), Alibaba Cloud\par
  \ifresearch
  - Developed multi-agent reinforcement learning for city-scale traffic signal control,
  reducing mean travel time by 2.8\% against the CoLight state of the art on real-city road
  networks of up to 196 intersections, and deployed it in production urban traffic signal
  control (IJCAI, IEEE TKDE, IJMLC).\par
  - Developed progressive transfer learning under distribution shift, raising Rank-1 accuracy
  by up to 4.5 points over the MGN backbone across four person re-identification benchmarks
  (MSMT17, Market-1501, CUHK03, DukeMTMC-reID) (IEEE TIP).
  \else
  - Developed and deployed multi-agent reinforcement learning for city-scale traffic signal
  control, cutting mean travel time against the prior state of the art on real road networks of
  up to 196 intersections (IJCAI, IEEE TKDE, IJMLC).\par
  - Developed progressive transfer learning under distribution shift for production vision
  models (IEEE TIP).
  \fi \\
\gap\textbf{Research Intern}\par 2018--2021 &
  \gap DAMO Academy, Alibaba Group\par
  - Developed large-scale multi-agent reinforcement learning and deep representation learning
  methods for the traffic engine of Alibaba City Brain; 7 granted national invention
  patents. \\
\end{entries}
}

\ifresearch
  \educationsection
  \employmentsection
\else
  \employmentsection
  \educationsection
\fi

% %========================================================== awards
% \section{Awards and Honours}

% \begin{datedentries}[3cm]
% \gap[4pt]\textbf{Alibaba Cloud Best Employee Award} & \gap[4pt]2024, 2025 \\
% \gap\textbf{Alibaba DAMO Academy Outstanding Intern Award} & \gap 2018, 2019, 2021 \\
% \gap\textbf{Zhejiang University Outstanding Graduate Student Award} & \gap 2019, 2020 \\
% \end{datedentries}

%========================================================== publications
\section{Selected Publications}

The $*$ symbol denotes equal contribution.

\subsection{Conference Papers}

\begin{publist}
\pub{C4} Weihang Pan, \me, Yuxiang Zhang, Zhongming Jin, Binbin Lin, Xiaofei He, Jieping Ye.
  ``ChainPrune: Evaluating and Reducing Redundancy in Long Chain-of-Thought Reasoning''
  \venue{ISWC}, 2026.

\pub{C3} Weihang Pan, \me, Yong Wu, Xun Liang, Zhongming Jin, Qiang Fu, Penghui Shang, Binbin Lin,
  Xiaofei He, Jieping Ye. ``FGD-Align: Pluralistic Alignment for Large Language Models via Fuzzy
  Group Decision-Making'' \venue{AAAI Conference on Artificial Intelligence (AAAI)}, 2026.

\pub{C2} Yuxiang Zhang*, \me*, Weihang Pan, Zhongming Jin, Qiang Fu, Deng Cai, Binbin Lin,
  Jieping Ye. ``TokenSqueeze: Performance-Preserving Compression for Reasoning LLMs''
  \venue{Neural Information Processing Systems (NeurIPS)}, 2025.

\pub{C1} \me*, Shuxian Liang*, Long Wei, Zhongming Jin, Jianqiang Huang, Deng Cai, Xiaofei He,
  Xian-Sheng Hua. ``MaCAR: Urban Traffic Light Control via Active Multi-agent Communication and
  Action Rectification'' \venue{International Joint Conference on Artificial Intelligence
  (IJCAI)}, 2020.
\end{publist}

\subsection{Journal Articles}

\begin{publist}

\ifresearch
\pub{J4} Weihang Pan, Binbin Lin, Yafei Wang, \me, Xinkui Zhao, Xiaofei He, Jieping Ye.
  ``Cooperative Driving at Multiple Unsignalized Intersections in Fully Autonomous Driving
  Scenarios'' \venue{IEEE Transactions on Intelligent Transportation Systems (T-ITS)},
  vol.~26, no.~12, 2025.

\pub{J3} Chao Xiang, Zhongming Jin, \me, Xian-Sheng Hua, Yao Hu, Wei Qian, Kaili Zhu, Deng Cai,
  Xiaofei He. ``Optimizing Traffic Efficiency via a Reinforcement Learning Approach Based on Time
  Allocation'' \venue{International Journal of Machine Learning and Cybernetics (IJMLC)},
  vol.~14, no.~10, 2023.
\fi

\pub{J2} Xin Guo, \me, Pengfei Wang, Zhongming Jin, Jianqiang Huang, Deng Cai, Xiaofei He,
  Xian-Sheng Hua. ``Urban Traffic Light Control via Active Multi-Agent Communication and
  Supply-Demand Modeling'' \venue{IEEE Transactions on Knowledge and Data Engineering (TKDE)},
  vol.~35, no.~4, 2023. (Journal extension of C1.)

\pub{J1} \me, Dong Shen, Zhongming Jin, Jianqiang Huang, Deng Cai, Xian-Sheng Hua.
  ``Progressive Transfer Learning'' \venue{IEEE Transactions on Image Processing (TIP)},
  vol.~31, 2022.
\end{publist}

\subsection{Technical Reports}

\begin{publist}
\pub{P1} \me, Yu Fu, Zhiyuan He, Yuxuan Huang, Ka Yiu Lee, Meng Fang, Weilin Luo, Jun Wang.
  ``OneManCompany: An Open-Source Operating System for Self-Organizing Multi-Agent Systems''
  \venue{arXiv preprint arXiv:2604.22446}, 2026.
  \href{https://arxiv.org/abs/2604.22446}{Paper}.
\end{publist}

% \subsection{Theses}

% \begin{publist}
% \pub{T1} \me. ``Stable Learning for Non-I.I.D Data'' Ph.D.\ thesis, Zhejiang University, 2021.
% \end{publist}

%========================================================== talks
% \section{Invited Talks}

%========================================================== service
% \section{Professional Service}

% \subsection{Conference Area Chair}

% \subsection{Conference Refereeing}

% Add years and switch to \begin{datedentries} (as used above) once known, e.g.
%   \gap Neural Information Processing Systems (NeurIPS) & \gap 2023--2026 \\
% \begin{plainentries}
% \item Neural Information Processing Systems (NeurIPS)
% \item International Conference on Learning Representations (ICLR)
% \item AAAI Conference on Artificial Intelligence (AAAI)
% \item International Joint Conference on Artificial Intelligence (IJCAI)
% \item European Conference on Computer Vision (ECCV)
% \item IEEE Transactions on Image Processing (TIP)
% \item IEEE Transactions on Multimedia (TMM)
% \item IEEE Transactions on Cognitive and Developmental Systems (TCDS)
% \end{plainentries}

% \subsection{Workshop Organization}

% \subsection{Seminar Organization}

\end{document}
